\chapter{绪论}\label{chap:introduction}

本章展示论文写作中常用的排版元素，包括标题层级、列表和字体等基础功能。在正式撰写毕业设计（论文）时，请将本章内容替换为实际的绪论内容，通常包括课题研究背景与意义、国内外研究现状综述、本文主要研究内容与组织结构等。

\section{标题层级示例}

本模板按照同济大学规范配置了三级标题层级。编号体系因专业类别而异，由 \verb|main.tex| 中的 \texttt{field} 选项控制，具体对应关系如\tabref{tab:heading-levels}所示。官方撰写规范对三级以下标题不作强制要求，本模板仍配置了 \cs{subsubsection} 与 \cs{paragraph} 两级作为参考样式，可按需自行使用。

\begin{table}[htbp]
  \centering
  \caption{标题层级与编号体系（\texttt{field} 选项）}\label{tab:heading-levels}
  \renewcommand{\arraystretch}{1.8}
  \begin{tabular}{cclll}
    \toprule
    \textbf{层级} & \textbf{\LaTeX{} 命令} & \textbf{理工科（默认）} & \textbf{文科} & \textbf{示例（理工科）} \\
    \midrule
    1 & \texttt{\textbackslash chapter}    & 1, 2, 3      & 一、二、三   & {\heiti\tjfontchapter 1\quad 一级标题示例} \\
    2 & \texttt{\textbackslash section}    & 1.1, 1.2     & （一）（二） & {\heiti\tjfontbody 1.1\hspace{0.5em}二级标题示例} \\
    3 & \texttt{\textbackslash subsection} & 1.1.1, 1.1.2 & 1.\ 2.       & {\heiti\tjfontbody 1.1.1\hspace{0.5em}三级标题示例} \\
    \bottomrule
  \end{tabular}
\end{table}

二级标题用于划分各章的主要内容板块。例如，在绪论中可使用二级标题分别介绍“研究背景与意义”、“国内外研究现状”和“本文组织结构”等。

\subsection{三级标题示例}

三级标题用于在二级标题下进一步细分内容，例如在“国内外研究现状”下分别介绍国内和国外的研究进展。在实际论文写作中，通常不建议使用超过三级的标题层级，以保持文章结构的清晰性。

\section{列表}

列表是论文中常用的排版元素，本节展示无序列表和有序列表的使用方法。有序列表的编号格式已按照同济大学毕业设计（论文）撰写规范进行配置。

\subsection{无序列表}

无序列表适用于列举不需要特定顺序的内容，使用 \texttt{itemize} 环境实现。

\begin{itemize}
    \item 同济大学概述：同济大学是一所以工科著称的综合性大学，具有深厚的历史和优秀的教学资源，是中国著名的高等学府之一，享有盛誉。
    \begin{itemize}
        \item 创建于1907年：同济大学始建于1907年，是中国历史最悠久的大学之一，拥有悠久的历史和丰富的文化底蕴；
        \item 位于上海市：校园位于中国上海市，这里是中国经济和文化的重要中心，为学生提供了广阔的发展空间，是学习和生活的理想之地，也是中国高等教育的重要基地；
        \item 以工科著称：同济大学在工程学科方面享有盛誉，是中国顶尖的工科学府之一，培养了大批优秀的工程技术人才，为国家的工程建设和科技创新做出了重要贡献。
    \end{itemize}
    \item 校园文化：同济大学注重校园文化建设，营造了丰富多彩的校园生活。这里的文化活动多样，包括学生社团、艺术节、体育赛事等，为学生提供了展示才华和发展兴趣的平台。
    \item 学术成就：同济大学在多个领域取得了显著的学术成就，培养了大批优秀人才。学校在科学研究、技术创新和社会服务方面都有突出的表现，为国家和社会的发展做出了重要贡献。
\end{itemize}

\subsection{有序列表}

有序列表适用于列举需要按顺序排列的内容，使用 \texttt{enumerate} 环境实现。

需要注意的是，根据同济大学提供的理工科毕业设计（论文）撰写规范，有序列表的第一级编号应使用全角圆括号内的数字，如“（1）”、“（2）”等；第二级编号应使用圆圈内的阿拉伯数字，如“\Circled{1}”、“\Circled{2}”等。此外，有序列表的第二级应为行内列表，即不应另起一行，而应与上一级列表项在同一行内；我们在写 \LaTeX{} 代码时，应该使用 \texttt{enumerate*} 环境来实现有序列表的第二级。

\begin{enumerate}
    \item 同济大学历史：同济大学是中国著名的高等学府之一，拥有悠久的历史和丰富的文化底蕴。以下是同济大学的重要历史节点：
    \begin{enumerate*}
        \item 1907年：创建，最初为德文医学堂；
        \item 1927年：成为国立大学，更名为同济大学；
        \item 1952年：调整学科设置，成为以工科为主的综合性大学。
    \end{enumerate*}
    \item 学术体系：同济大学在学术方面具有很高的声誉，涵盖多个重要学科领域。主要学科包括：
    \begin{enumerate*}
        \item 工程学，包括土木工程、建筑工程、机械工程等；
        \item 建筑学，包括建筑设计、城市规划等；
        \item 医学，包括临床医学、药学等。
    \end{enumerate*}
    \item 国际合作：同济大学积极参与国际交流与合作，拓展全球视野，提升国际影响力。主要的合作项目包括：
    \begin{enumerate*}
        \item 德国交流项目，包括学生交换、教师访问等；
        \item 全球合作伙伴，包括国际知名大学、科研机构等；
        \item 国际学术会议，包括学术研讨会、国际学术交流等。
    \end{enumerate*}
\end{enumerate}

\section{字体}

本模板正文默认使用宋体，标题使用黑体。以下展示模板中可用的中文字体及其效果：

\begin{itemize}
    \item {\songti 宋体}：使用命令 \texttt{\textbackslash songti}。
    \item {\heiti 黑体}：使用命令 \texttt{\textbackslash heiti}。
    \item {\fangsong 仿宋}：使用命令 \texttt{\textbackslash fangsong}。
    \item {\kaishu 楷书}：使用命令 \texttt{\textbackslash kaishu}。
\end{itemize}

\subsection{字号设置}

本模板已按照同济大学毕业设计（论文）撰写规范预设了各部分的字号，一般无需手动调整。如有特殊需要，可使用 \verb|\zihao{}| 命令设置字号，例如 \verb|\zihao{-4}| 表示小四号。常用字号对应关系如\tabref{tab:font-sizes}所示。

\begin{table}[htbp]
  \centering
  \caption{常用字号对照}\label{tab:font-sizes}
  \begin{tabular}{ccl}
    \toprule
    \textbf{字号命令} & \textbf{字号名称} & \textbf{用途} \\
    \midrule
    \texttt{\textbackslash zihao\{-2\}} & 小二号 & 课题名称、摘要页标题 \\
    \texttt{\textbackslash zihao\{-3\}} & 小三号 & 章标题、封面字段值 \\
    \texttt{\textbackslash zihao\{3\}}  & 三号   & 摘要副标题（\texttt{\textbackslash tjfontinfotitle}：信息说明页标题） \\
    \texttt{\textbackslash zihao\{4\}}  & 四号   & 摘要、目录、参考文献、谢辞标题 \\
    \texttt{\textbackslash zihao\{-4\}} & 小四号 & 正文、节标题、页眉页脚、目录内容、公式 \\
    \texttt{\textbackslash zihao\{5\}}  & 五号   & 表题、图题、装订线文字 \\
    \texttt{\textbackslash zihao\{-5\}} & 小五号 & 表内文字 \\
    \bottomrule
  \end{tabular}
\end{table}

\subsection{文字样式}

常用的文字样式命令包括：\textbf{加粗}（\verb|\textbf{}|）、\textit{斜体}（\verb|\textit{}|）、\underline{下划线}（\verb|\underline{}|）以及\emph{强调}（\verb|\emph{}|）。在中文环境下，\verb|\emph{}| 会自动切换为楷书以示强调。

{\songti

【宋体示例】1900年前后，由埃里希·宝隆创办的“同济医院”正式挂牌。医院的医师大都是“德医公会”成员。他们白天忙于经营自己的诊所，只有傍晚到医院看门诊、动手术。埃里希·宝隆医生看到医院里的医疗力量不足，计划在院内设立一所德文医学堂，招收中国学生，以培养施诊医生。这个计划得到德国驻沪总领事以及德国政府高等教育司的支持。1906年，他们设立了一个支持医学堂开办的基金会，得到了德国“促进德国与外国思想交流的科佩尔基金会”的协助，筹集到一批医科书刊及新式的外科手术电动器械等物品。}

{\heiti 【黑体示例】1907年6月医学堂开学前，德国驻沪总领事克纳佩在上海不仅号召德国商人捐款，而且要求德国洋行向中国商人募捐。同时，费舍尔还要求中国官方的资助和支持，克纳佩利用在中德两国募来的捐款，成立了“为中国人办的德国医学堂基金会”。当时规定，捐款金额较多者可成为医学堂董事会董事。医学堂建立时定名为德文医学堂，并成立了董事会负责学校的管理。董事会由18人组成，主要成员有：三个德医公会元老：宝隆、福沙伯（第二任校长）、福尔克尔；三名德国商人：莱姆克、米歇劳和赖纳；两名中国绅商：朱葆三（沪军都督府财政部长及上海商务会会长，大买办）、虞洽卿（荷兰银行买办）；总领事馆的副领事弗赖海尔·冯·吕特等。埃里希·宝隆医生被正式推选为董事会总监督（董事长）兼学堂首任总理（校长），负责学堂的管理。医学堂的校址设在同济医院对面的白克路（今凤阳路415号上海长征医院内）。1907年10月1日德文医学堂举行了开学典礼。}

{\ifcsname fangsong\endcsname\fangsong【仿宋示例】\else[无 \cs{fangsong} 字体。]\fi
1923年3月17日北洋政府教育部下达第108号训令，批准同济工科“改为大学”。学校随即召开董事会议，将学校定名为“同济大学”。1923年3月26日，学校以“同济大学董事会”名义呈文北洋政府教育部，称“经校董会议定名称为同济大学”。1923年4月24日，北洋政府教育部下达第634号“指令”，称“该校名称拟改为同济大学，应予照准备案”。1924年5月20日北洋政府教育部下达第120号训令，批准同济医科为大学。从此以后，5月20日定为校庆日。}

{\ifcsname kaishu\endcsname\kaishu【楷书示例】\else[无 \cs{kaishu} 字体。]\fi
抗日战争胜利后，1946年，国立同济大学分批迁回上海。由于缺少校舍，学校分散教学，成为斜跨上海市区、分散十多处的“大学校”。其中学校办公室和医学院位于善钟路100弄10号（今常熟路），附属医院分别为白克路上的中美医院（原宝隆医院，今凤阳路）和同孚路82号（今石门一路）的原德国医院，理学院在平昌街日本第七国民学校内（今国顺路上海电视大学），文法学院在四川北路（今复兴初级中学），新生院位于江湾新市区的市图书馆（今黑山路），高级工业职业学校则在江湾魏德迈路（今邯郸路），附属中学位于市博物馆（今长海医院飞机楼）。工学院位于其美路的原日本中学（今杨浦区四平路1239号），1949年后逐步发展成同济大学的主校区。}

